%% -------------------------------------
%%
%% Configuration file.
%%
%% ---------------------
%% Change log:
%% 	2026-03-15 - initial version.
%%
%% ---------------------
%% Authors:
%% 	Santiago Husain Cerruti - hola dot santiago at proton dot me
%%
%% ---------------------
%% Filename: config.tex
%% Part of FIUBA-thesis project.
%% Copyright (c) 2026 Santiago Husain Cerruti.
%% Licensed under the MIT License.
%% See LICENSE file for full text.
% -------------------------------------
% -------------------------------------------------------
% FORMAT
% -------------------------------------------------------
% Input encoding and language.
\usepackage[spanish, es-tabla]{babel}
\usepackage[utf8]{inputenc}
% Refer to p.92 of "The not so short introduction to Latex" 
% by Tobias Oetiker to understand what role have this packages.
\usepackage{lmodern}
\usepackage[T1]{fontenc}
\usepackage{textcomp}
% Set Arial font.
\usepackage[scaled]{helvet} 
\renewcommand*\familydefault{\sfdefault}

% The mitthesis template uses Letter-sized pages. We must adapt it to A4 size.
\usepackage[a4paper, left=3.5cm, right=2.5cm, top=3cm, bottom=3cm]{geometry}

% Provide a high-level, flexible interface for creating document-level commands and environments.
\usepackage{xparse}

% Makes the List of Tables (LoT) and List of Figures (LoF) appear in the Table of Contents (ToC).
% The 'nottoc' option makes the Table of Contents itself not to be listed in the ToC.
\usepackage[nottoc]{tocbibind}

% Increase number-column width in LoF/LoT to avoid overfull boxes with multi-level numbering.
%Your figure numbering like 1.1.1 was wider than the default number slot in LoF.
%Widening that slot prevents line-3 overflow in FIUBA-thesis-main.lof.
\makeatletter
\renewcommand*\l@figure{\@dottedtocline{1}{1.5em}{3.2em}}
\renewcommand*\l@table{\@dottedtocline{1}{1.5em}{3.2em}}
\makeatother

% ------------------------------------------------------------------------
% PAGE MARGINS
% ------------------------------------------------------------------------
% Change this lengths as needed. Usually it's not necessary.
% Definitions:
%	\textwidth	 	the width of the text on the page.
%	\hoffset 		the horizontal offset of the text (ie: how far the text is from the left of the page).
%	\textheight 	the height of the text on the page.
%	\voffset 		the vertical offset of the text (ie: how far the text is from the top of the page).
%
% Modify/comment this parameters according to the user preferences. 
% Adds cm to the default width (if used with hoffset, centering is not altered)
%\addtolength{\textwidth}{4cm}
% Removes cm to the margin, but keeping the aspect relation of each margin.
%\addtolength{\hoffset}{-2cm}
%\addtolength{\textheight}{4cm}
%\addtolength{\voffset}{2cm}
\setlength{\headheight}{15.1pt}
% ------------------------------------------------------------------------


% To add pictures to a page at absolute positions. Useful to set the cover image.
\usepackage{eso-pic}

% Verbatim text formatting, including headers.
\usepackage{fancyvrb}
% Boxes management.
\usepackage{fancybox}
% Headers management.
\usepackage{fancyhdr}
	% First you need to clear the default layout.
	\fancyhead{}
	\fancyfoot{}
% -------------
% Uncomment this line when using two-side mode document.
% -------------
	% \fancyhead[RE]{\rightmark}
	\fancyhead[LO]{\rightmark}
	\fancyfoot[C]{\thepage}
	% Remove decorative lines.
	\renewcommand{\headrulewidth}{0pt}
	\renewcommand{\footrulewidth}{0pt}
	\pagestyle{fancy}

% To remove the page numbers uncomment the following two lines:
%\thispagestyle{empty}
%\setcounter{savepage}{\thepage}

% Package for underlining, adding color to text, and highlighting text in color.
% Text highlighting package. It solves the 'fullbox' errors when using \colorbox{declared-color}{text}.
% It doesn't support letters with accents, so they must be included using the '\'' syntax.
% Usage: \hl{text}.
\usepackage{soul}
% Sets the 'soul' package to higlight text using color red.
\sethlcolor{red}

% -------------------------------------
% Easy quoting.
\usepackage[style=american]{csquotes}


% ------------------------------------------------------------------------
% MATHS
% ------------------------------------------------------------------------
\usepackage{amsmath,amsfonts,amsthm,amssymb,amsbsy}
% Equation numbering format by section. E.g., section x, equation y: (x.y)
\numberwithin{equation}{section}
\numberwithin{figure}{section}
\usepackage{mathtools}
% Useful for arrays of equations.
\usepackage[retainorgcmds]{IEEEtrantools}

% ------------------------------------------
% Setting Up Automatic Scaling Brackets.
% ------------------------------------------
% When you write equations with tall elements like fractions or matrices, standard parentheses ( and ) stay small and look awkward. While typing \left( and \right) fixes this, doing it every single time gets exhausting and clutters your code.
% Define custom paired delimiters
\DeclarePairedDelimiter\paren{(}{)}       % Parentheses
\DeclarePairedDelimiter\sqbrac{[}{]}      % Square brackets
\DeclarePairedDelimiter\curly{\{}{\}}     % Curly braces
\DeclarePairedDelimiter\abs{\lvert}{\rvert} % Absolute value
\DeclarePairedDelimiter\norm{\lVert}{\rVert} % Norm

% ----------------------------
% Imaginary unit.
% The international standard for mathematical typesetting (ISO 80000-2) dictates that mathematical constants—like the base of the natural logarithm ($e$), pi ($\pi$), and the imaginary unit ($i$ or $j$)—should be typeset in an upright (Roman) font to clearly distinguish them from variables. (Notice how the 'i' stands straight up when using it, making it instantly recognizable as a constant rather than a variable exponent).
% ----------------------------
\newcommand{\iu}{\mathrm{i}} % Imaginary unit for math/physics
\newcommand{\ju}{\mathrm{j}} % Imaginary unit for engineering
% ----------------------------
% Generic counter
% ----------------------------
\newcounter{hipotesis_ct}
\setcounter{hipotesis_ct}{1}

% ----------------------------
% Derivatives operators.
\usepackage{derivative}

% ----------------------------
% Set operators.
% ----------------------------
% We create a dummy command \given that we will redefine inside our \set command
\providecommand\given{} 
% The 'X' allows us to inject formatting inside the delimiters
\DeclarePairedDelimiterX\set[1]\{\}{
	% This makes \given print a vertical bar that scales with the outer braces
	\renewcommand\given{\;\delimsize\vert\;}
	#1
}

% ----------------------------
% Commands to write vectors.
% ----------------------------
\usepackage[c]{esvect}
\newcommand{\vecbf}[1]{\mathbf{#1}}


% ----------------------------
% Commands to write vectors.
% ----------------------------
\usepackage[c]{esvect} % Keep this ONLY if you actually use \vv{x} for arrows somewhere
\usepackage{bm}        % The modern bold math package

% Core semantic definitions
\newcommand{\mVec}[1]{\bm{#1}}       % Bold italic vectors (works on Greek too!)
\newcommand{\Mat}[1]{\bm{#1}}        % Matrices (often same as vectors, but good to separate semantically)
\newcommand{\quat}[1]{\mathbf{#1}}   % Quaternions (Standard is upright bold)

% Note: The \dot goes OUTSIDE the \mVec so the dots don't get bolded.
% Vector decorations
\newcommand{\vechat}[1]{\hat{\mVec{#1}}}
\newcommand{\vecunit}[1]{\breve{\mVec{#1}}}
\newcommand{\vecdot}[1]{\dot{\mVec{#1}}}      % Velocity
\newcommand{\vecddot}[1]{\ddot{\mVec{#1}}}    % Acceleration
\newcommand{\vecdddot}[1]{\dddot{\mVec{#1}}}  % Jerk
\newcommand{\vecprod}[2]{#1 \times #2}

% Quaternion operations
\newcommand{\quatComp}[2]{#1 \circ #2}	% Composition.
\newcommand{\quatConj}[1]{#1^{\ast}}	% Conjugate.
\newcommand{\quatRot}[2]{\Re(#1, #2)} 	% Rotation.
% ----------------------------
% Estimators.
% Dynamically stretches the hat to perfectly cover whatever is inside of it.
\newcommand{\estimated}[1]{\widehat{#1}}

% ----------------------------
% Command for equations definitions.
% Requires the amssymb package
\newcommand{\defineEq}{\triangleq}

% ----------------------------
% Système international d'unités
\usepackage{siunitx}

% ---------------------------------
% Automatic matrix transposition.
% ---------------------------------
% This snippet uses LaTeX3 (expl3) programming to create a highly useful new math environment called bmatrixT. In short, it automatically transposes a matrix. You type the matrix normally, and when the document compiles, the code swaps the rows and columns and wraps the result in standard square brackets (a bmatrix).
% It also works with pmatrix (parentheses) and vmatrix (vertical bars).
%
% How it works (using the bmatrixT environment as example):
%
% Capturing the Content: It uses the environ package to grab everything you type inside \begin{bmatrixT} ... \end{bmatrixT} and saves it as a single chunk of text (\BODY).
%
% Deconstructing the Matrix: It splits the input at every line break (\\) to figure out the rows.
%
% It then splits each row at every ampersand (&) to separate the individual column items.
%
% It stores every single item in a temporary property list using its coordinates (Row, Column) as a tag.
%
% Rebuilding the Matrix (The Transpose): It uses nested loops to rebuild the matrix, but it reverses the order. It loops through the columns first, and the rows second.
%
% This effectively swaps the X and Y coordinates of every item.
%
% Formatting the Output: Finally, it takes the newly swapped items, adds the & and \\ symbols back in the correct spots, and feeds the whole thing into a standard bmatrix environment.
\usepackage{environ}

\ExplSyntaxOn
% Declare variables
\int_new:N \l_marine_transpose_row_int
\int_new:N \l_marine_transpose_col_int
\seq_new:N \l_marine_transpose_rows_seq
\seq_new:N \l_marine_transpose_arow_seq
\prop_new:N \l_marine_transpose_matrix_prop
\tl_new:N \l_marine_transpose_last_tl
\tl_new:N \l_marine_transpose_body_tl

% 1. The function takes TWO arguments (:nn)
% #1 is the environment type (bmatrix, pmatrix, etc.)
% #2 is the matrix content
\cs_new_protected:Nn \marine_transpose:nn
{
	\seq_set_split:Nnn \l_marine_transpose_rows_seq { \\ } { #2 }
	\int_zero:N \l_marine_transpose_row_int
	\prop_clear:N \l_marine_transpose_matrix_prop
	\seq_map_inline:Nn \l_marine_transpose_rows_seq
	{
		\int_incr:N \l_marine_transpose_row_int
		\int_zero:N \l_marine_transpose_col_int
		\seq_set_split:Nnn \l_marine_transpose_arow_seq { & } { ##1 }
		\seq_map_inline:Nn \l_marine_transpose_arow_seq
		{
			\int_incr:N \l_marine_transpose_col_int
			\prop_put:Nxn \l_marine_transpose_matrix_prop
			{
				\int_to_arabic:n { \l_marine_transpose_row_int }
				,
				\int_to_arabic:n { \l_marine_transpose_col_int }
			}
			{ ####1 }
		}
	}
	\tl_clear:N \l_marine_transpose_body_tl
	\int_step_inline:nnnn { 1 } { 1 } { \l_marine_transpose_col_int }
	{
		\int_step_inline:nnnn { 1 } { 1 } { \l_marine_transpose_row_int }
		{
			\tl_put_right:Nx \l_marine_transpose_body_tl
			{
				\prop_item:Nn \l_marine_transpose_matrix_prop { ####1,##1 }
				\int_compare:nF { ####1 = \l_marine_transpose_row_int } { & }
			}
		}
		\tl_put_right:Nn \l_marine_transpose_body_tl { \\ }
	}
	
	% 2. Wrap the output in the environment specified by #1
	\begin{#1}
		\l_marine_transpose_body_tl
	\end{#1}
}

% 3. Generate a variant that expands the second argument (nV)
\cs_generate_variant:Nn \marine_transpose:nn { nV }
\cs_generate_variant:Nn \prop_put:Nnn { Nx }

% 4. Define all the specific transpose environments
\NewEnviron{bmatrixT}{ \marine_transpose:nV { bmatrix } \BODY }
\NewEnviron{pmatrixT}{ \marine_transpose:nV { pmatrix } \BODY }
\NewEnviron{vmatrixT}{ \marine_transpose:nV { vmatrix } \BODY }
\NewEnviron{BmatrixT}{ \marine_transpose:nV { Bmatrix } \BODY } % Braces {}
\NewEnviron{VmatrixT}{ \marine_transpose:nV { Vmatrix } \BODY } % Double bars ||

\ExplSyntaxOff

% ---------------------------------------------------
% GRAPHICS AND TABLES.
% ---------------------------------------------------
% Package for improving the interface with floating objects, such as graphics.
\usepackage{float}
\usepackage{graphicx}
% Packages that offer greater flexibility when using captions for figures, etc.
\usepackage{caption}
\usepackage{subcaption}
\usepackage{epstopdf}
% -----------------------
\usepackage{booktabs}
\usepackage{multirow}

% ---------------------------------------------------
% BIBLIOGRAPHY AND APPENDICES.
% ---------------------------------------------------
% The Backend: biblatex works best with a sorting engine called biber. Explicitly stating backend=biber prevents LaTeX from accidentally trying to use older, broken engines.
\usepackage[
	backend=biber,
	style=ieee,       	% Pick your citation style (numeric, apa, ieee, authoryear)
	maxbibnames=99,   	% Usually, you want all authors in the final list...
	minbibnames=1,  	% Ensures no conflict in the bibliography
	maxcitenames=2,   	% ...but only 1 or 2 authors when citing inline
	mincitenames=1,
	uniquelist=false
]{biblatex}

% Use 'family-given' instead of 'last-first' which is technically deprecated in modern LaTeX to avoid hidden warnings.
\DeclareNameAlias{sortname}{family-given}
\DeclareNameAlias{default}{family-given}

% \adddot ensures the period is handled correctly.
\DefineBibliographyStrings{spanish}{%
	andothers = {et al\adddot},
}

% The .bib file to be used
\addbibresource{back-matter-bibliography.bib}
% Add a bibliography-only line-breaking relaxation to avoid overflow caused by a long bibliography line.
% \sloppy at bibliography start gives TeX more flexibility to break those long entries cleanly.
\AtBeginBibliography{\sloppy}

% ---------------------------------------------------
% GLOSSARY / ACRONYMS
% ---------------------------------------------------

% --------------------------------
% Spanish localization bug fix.
% Glossaries pulls in datatool, and the TeX setup only has English datatool locale files installed. With babel set to spanish, tracklang warns that no Spanish datatool locale exists. Disabling datatool locale warnings suppresses this specific non-fatal warning cleanly.
% Currently can’t install an official Spanish datatool locale, because TeX Live doesn’t ship one. datatool provides base locale files plus datatool-english and some region files, but no datatool-spanish file. The warning comes from datatool locale detection (via tracklang) when using babel with spanish.
% Alternative solution: create a custom local Spanish datatool .ldf in the project so we can remove lang-warn=false.
% Loading datatool-base with lang-warn=false disables Spanish locale warnings before glossaries triggers them.
% Use this: \usepackage[lang-warn=false]{datatool-base}
% Or this: \PassOptionsToPackage{lang-warn=false}{datatool-base}
% --------------------------------
% Required for the aligned glossary style
\usepackage{longtable}
% The 'acronym' option creates a dedicated list. The 'toc' option includes the glossary in the TOC.
\usepackage[toc,acronym,nomain,nonumberlist]{glossaries}
% No external programs needed to sort the entries. 
\makenoidxglossaries

% Force the acronyms to be bold.
\renewcommand{\glsnamefont}[1]{\textbf{#1}}

% The acronyms definitions will be aligned, and not printed like a standard dictionary (jagged text).
% Create a Custom "Flush Left" Style.
% Tweak the text width parameter 'p{0.XX\textwidth}' accordingly. It tells the second column (the acronyms descriptions) to take up XX% of the page width and wrap nicely like a normal paragraph.
% 1. Define your custom strictly left-justified style
\newglossarystyle{flushleft}{%
	\setglossarystyle{long}% Base it on the existing 'long' style
	\renewenvironment{theglossary}%
	{\begin{longtable}{@{} l p{0.85\textwidth} @{}}}
		{\end{longtable}}%
}
% 2. Tell the package to use the new style
\setglossarystyle{flushleft}

% Load the definitions here in the preamble (not in the document).
%% -------------------------------------
%%
%% Acronyms page file.
%%
%% ---------------------
%% Change log:
%% 	2026-03-15 - initial version.
%%
%% ---------------------
%% Authors:
%% 	Santiago Husain Cerruti - hola dot santiago at proton dot me
%%
%% ---------------------
%% Filename: front-matter-glossary.tex
%% Part of FIUBA-thesis project.
%% Copyright (c) 2026 Santiago Husain Cerruti.
%% Licensed under the MIT License.
%% See LICENSE file for full text.
% -------------------------------------
\newacronym{utc}{UTC}{Coordinated Universal Time}
\newacronym{adt}{ADT}{Atlantic Daylight Time}
\newacronym{est}{EST}{Eastern Standard Time}
\newacronym{SLR}{SLR}{Satellite Laser Ranging}
\newacronym{GPS}{GPS}{Global Positioning System}
\newacronym{IAU}{IAU}{International Astronomical Union}
\newacronym{FK5}{FK5}{Fundamental Katalog}
\newacronym{ECI}{ECI}{Earth-centered inertial}
\newacronym{ITRF}{ITRF}{International Terrestrial Reference Frame}
\newacronym{ECEF}{ECEF}{Earth-centered Earth-fixed}
\newacronym{IERS}{IERS}{International Earth Rotatation Service}
\newacronym{IRM}{IRM}{IERS Reference Meridian}
\newacronym{IRP}{IRP}{International Reference Pole}
\newacronym{CEP}{CEP}{Celestial Ephemeris Pole}
\newacronym{MOD}{MOD}{Mean Equinox of Date}
\newacronym{TOD}{TOD}{True Equator of Date}
\newacronym{GAST}{GAST}{Greenwich Apparent Sidereal Time}
\newacronym{MST}{MST}{Mean Solar Time}
\newacronym{UT}{UT}{Universal Time}
\newacronym{NASA}{NASA}{National Aeronautics and Space Administration}
\newacronym{JPL}{JPL}{Jet Propulsion Laboratory}
\newacronym{SFU}{SFU}{Solar Flux Unit}
\newacronym{NOAA}{NOAA}{National Oceanic and Atmospheric Administration}
\newacronym{EAM}{EAM}{Exponential Atmospheric Model}
\newacronym{SAM}{SAM}{Standard Atmosphere Model}
\newacronym{HP}{HP}{Harris-Priester}
\newacronym{SRP}{SRP}{Solar Radiation Pressure}
\newacronym{LOS}{LOS}{Line-of-Sight}
\newacronym{BC}{BC}{Ballistic Coefficient}
\newacronym{FEUV}{F \textsubscript{EUV}}{Extreme Ultraviolet Radiation}
\newacronym{EGM-96}{EGM-96}{Earth Gravity Model 1996}
\newacronym{TDRSS}{TDRSS}{NASA's Tracking and Data Relay Satellite System}
\newacronym{COTS}{COTS}{Commercial Off-The-Shelf}
\newacronym{TG}{TG}{Trajectory generator}
\newacronym{GT}{GT}{Generador de trayectorias}
\newacronym{CONAE}{CONAE}{Comisión Nacional de Actividades Espaciales}
\newacronym{POE}{POE}{Precise Orbit Ephemeris}
\newacronym{STK}{STK}{Systems Tool Kit}
\newacronym{HPOP}{HPOP}{High Precision Orbit Propagator}
\newacronym{AROM}{AROM}{Autonomous Reconfiguration and Orbit Maintenance}
\newacronym{SAR}{SAR}{Synthetic-Aperture Radar}
\newacronym{MOE}{MOE}{Mean Orbital Elements}
\newacronym{ROE}{ROE}{Relative Orbital Elements}
\newacronym{OE}{OE}{Orbital Elements}
\newacronym{SAOCOM}{SAOCOM}{Satélite Argentino de Observación Con Microondas}
\newacronym{LEO}{LEO}{Low-Earth Orbit}
\newacronym{DSS}{DSS}{Distributed Satellite Systems}
\newacronym{GNC}{GNC}{Guiado, Navegación y Control}
\newacronym{TCP}{TCP}{Transmission Control Protocol}
\newacronym{UDP}{UDP}{User Datagram Protocol}
\newacronym{XML}{XML}{Extensible Markup Language}
\newacronym{CSV}{CSV}{Comma-Separated Values}
\newacronym{YAML}{YAML}{YAML Ain't Markup Language}
\newacronym{RF}{RF}{Radiofrequencia}
\newacronym{CAN}{CAN}{Controller Area Network}
\newacronym{MVC}{MVC}{Model-View-Controller}
\newacronym{OP}{OP}{Observer Pattern}
\newacronym{ENP}{ENP}{Event Notification Pattern}

















	
% ---------------------------------------------------
% MISC.
% ---------------------------------------------------
% This pack­age sim­pli­fies the in­clu­sion of ex­ter­nal multi-page PDF doc­u­ments in LaTeX doc­u­ments. 
\usepackage{pdfpages}
\setboolean{@twoside}{false}

% ------------------------------------------------------------------------
% SOURCE CODE
% ------------------------------------------------------------------------
\usepackage[dvipsnames,svgnames]{xcolor} 
\usepackage{listings}
\usepackage{caption} % Better handling of code captions
% Algorithms and pseudocode handling.
\usepackage{algorithm}
\usepackage{algpseudocode}
\makeatletter
\renewcommand{\ALG@name}{Algoritmo}
\renewcommand{\listalgorithmname}{Lista de \ALG@name s}
\algrenewcommand\algorithmicensure{\textbf{Salida:}}
\algrenewcommand\algorithmicrequire{\textbf{Entrada:}}
\algnewcommand{\LineComment}[1]{\State \(\triangleright\) #1}
\makeatother

% ------------------------------------------------------------------------
% GLOBAL SETTINGS (Applied to all code)
% ------------------------------------------------------------------------
\lstset{
	inputencoding=utf8,
	extendedchars=true,
	literate={á}{{\'a}}1 {é}{{\'e}}1 {í}{{\'i}}1 {ó}{{\'o}}1 {ú}{{\'u}}1 {ñ}{{\~n}}1 {Á}{{\'A}}1 {É}{{\'E}}1 {Í}{{\'I}}1 {Ó}{{\'O}}1 {Ú}{{\'U}}1 {ü}{{\"u}}1 {Ü}{{\"U}}1,
	basicstyle=\ttfamily\small,
	breaklines=true,			% Enable wrapping
	breakatwhitespace=false,    % Wrap even if there isn't a space
	captionpos=b,
	keepspaces=true,
	showstringspaces=false,
	upquote=true,
	columns=fullflexible,       % Improves character spacing after wrapping
	postbreak=\mbox{\textcolor{red}{$\hookrightarrow$}\space} % Adds a little red arrow at the wrap point.
}

% Spanish Localization
\renewcommand{\lstlistingname}{Código fuente}       	% For the caption
\renewcommand{\lstlistlistingname}{Índice de códigos} 	% For the table of contents

% ------------------------------------------------------------------------
% DEFINE CUSTOM STYLES
% ------------------------------------------------------------------------

% 1. C/C++ Style
\lstdefinestyle{StyleC}{
	language=C++,
	commentstyle=\color{ForestGreen},
	keywordstyle=\color{RoyalBlue},
	stringstyle=\color{BrickRed},
	numberstyle=\tiny\color{gray},
	numbers=left,
	frame=lines,
	backgroundcolor=\color{gray!5}
}

% 2. YAML Style
\lstdefinelanguage{yaml}{
	keywords={true,false,null,y,n},
	keywordstyle=\color{darkgray}\bfseries,
	sensitive=false,
	comment=[l]{\#},
	morecomment=[s]{/*}{*/},
	commentstyle=\color{purple}\ttfamily,
	stringstyle=\color{blue}\ttfamily,
	moredelim=[l][\color{orange}]{\&},
	moredelim=[l][\color{magenta}]{*},
	moredelim=**[il][\color{red}{:}\color{blue}]{:},
	morestring=[b]',
	morestring=[b]",
	frame=lines,
	backgroundcolor=\color{gray!5}
}

% 3. Terminal Style (For console output)
\lstdefinestyle{terminal}{
	backgroundcolor=\color{black},
	basicstyle=\ttfamily\small\color{white},
	frame=lines,
	numbers=none
}
